% GENERATED FILE. Edit canonical lessons, then run npm run generate:books.
% canonical-source-hash: fnv1a64:f37180786747fc84
% canonical-lessons: RU-C230-sadovnik, RU-C230-kassir, RU-C230-parikmakher, RU-C230-okhrannik, RU-C230-nyanya

\chapter{Gardener, Cashier, Hairdresser, Security guard, Nanny}
\label{ch:ru-gardener-cashier-hairdresser-security-guard-nanny}
\hlchaptermodality{230}

\begin{chapteropening}
\textbf{By the end of this chapter:} \emph{I can say a gardener, a cashier, a hairdresser, a security guard and a nanny.}

Complete the last lesson of chapter 230: I can say a gardener, a cashier, a hairdresser, a security guard and a nanny.
\end{chapteropening}

\section[\texorpdfstring{\ru{садовник} (sadóvnik)}{садовник (sadóvnik)}]{\ru{садовник} (sadóvnik) — a gardener}

\label{lesson:RU-C230-sadovnik}



\begin{quote}
Before the new one: say the Russian for a programmer, then the Russian for a lecturer.
\end{quote}

\subsection*{You'll want to know: \ru{садовник}}
\textbf{\ru{садовник}} (\emph{sadóvnik}) — \textquotedblleft{}a gardener\textquotedblright{}.

\begin{cousinweb}
\textbf{\ru{садовник}} is \emph{\ru{сад}}, \textquotedblleft{}a garden\textquotedblright{}, with \emph{-\ru{овник}}, \textquotedblleft{}a person\textquotedblright{}: a gardener. \emph{\ru{сад}} comes from \emph{\ru{садить}}, \textquotedblleft{}to plant\textquotedblright{}, the \textquotedblleft{}make sit\textquotedblright{} form of \emph{\ru{сидеть}}, \textquotedblleft{}to sit\textquotedblright{} (Vasmer).
\end{cousinweb}

\begin{grammarlens}[title={What you've built}]
One more word for this chapter.
\end{grammarlens}

\subsection*{Your turn}
\begin{itemize}
\raggedright
  \item \emph{Say it:} \emph{sadóvnik}
  \item \emph{Say it:} \emph{sadóvnik}, once more
  \item \emph{Say it:} say \emph{sadóvnik}
  \item \emph{Recall:} say the Russian for a programmer, then the Russian for a lecturer, then say \emph{sadóvnik} again
\end{itemize}

\begin{tcolorbox}[breakable,colback=teal!4,colframe=teal!35!black,title={Before you move on}]
What does \ru{садовник} mean? (\textquotedblleft{}A gardener\textquotedblright{}.) Say it once more.
\end{tcolorbox}
\section[\texorpdfstring{\ru{кассир} (kassír)}{кассир (kassír)}]{\ru{кассир} (kassír) — a cashier}

\label{lesson:RU-C230-kassir}



\begin{quote}
Before the new one: say the Russian for a lecturer, then the Russian for a gardener.
\end{quote}

\subsection*{You'll want to know: \ru{кассир}}
\textbf{\ru{кассир}} (\emph{kassír}) — \textquotedblleft{}a cashier\textquotedblright{}.

\begin{cousinweb}
\textbf{\ru{кассир}} — older \emph{\ru{касирер}}, from Peter the Great's time; it comes from German \emph{Kassierer}, from \emph{Kasse}, \textquotedblleft{}a cash box\textquotedblright{} (Vasmer).
\end{cousinweb}

\begin{grammarlens}[title={What you've built}]
One more word for this chapter.
\end{grammarlens}

\subsection*{Your turn}
\begin{itemize}
\raggedright
  \item \emph{Say it:} \emph{kassír}
  \item \emph{Say it:} \emph{kassír}, once more
  \item \emph{Say it:} say \emph{kassír}
  \item \emph{Recall:} say the Russian for a lecturer, then the Russian for a gardener, then say \emph{kassír} again
\end{itemize}

\begin{tcolorbox}[breakable,colback=teal!4,colframe=teal!35!black,title={Before you move on}]
What does \ru{кассир} mean? (\textquotedblleft{}A cashier\textquotedblright{}.) Say it once more.
\end{tcolorbox}
\section[\texorpdfstring{\ru{парикмахер} (parikmákher)}{парикмахер (parikmákher)}]{\ru{парикмахер} (parikmákher) — a hairdresser}

\label{lesson:RU-C230-parikmakher}



\begin{quote}
Before the new one: say the Russian for a gardener, then the Russian for a cashier.
\end{quote}

\subsection*{You'll want to know: \ru{парикмахер}}
\textbf{\ru{парикмахер}} (\emph{parikmákher}) — \textquotedblleft{}a hairdresser\textquotedblright{}.

\begin{cousinweb}
\textbf{\ru{парикмахер}} comes from German \emph{Perückenmacher}, \textquotedblleft{}wig-maker\textquotedblright{}, reshaped after \emph{\ru{парик}}, \textquotedblleft{}a wig\textquotedblright{}; older \emph{\ru{парукмакар}} came from Dutch \emph{paruikmaker} (Vasmer).
\end{cousinweb}

\begin{grammarlens}[title={What you've built}]
One more word for this chapter.
\end{grammarlens}

\subsection*{Your turn}
\begin{itemize}
\raggedright
  \item \emph{Say it:} \emph{parikmákher}
  \item \emph{Say it:} \emph{parikmákher}, once more
  \item \emph{Say it:} say \emph{parikmákher}
  \item \emph{Recall:} say the Russian for a gardener, then the Russian for a cashier, then say \emph{parikmákher} again
\end{itemize}

\begin{tcolorbox}[breakable,colback=teal!4,colframe=teal!35!black,title={Before you move on}]
What does \ru{парикмахер} mean? (\textquotedblleft{}A hairdresser\textquotedblright{}.) Say it once more.
\end{tcolorbox}
\section[\texorpdfstring{\ru{охранник} (okhránnik)}{охранник (okhránnik)}]{\ru{охранник} (okhránnik) — a security guard}

\label{lesson:RU-C230-okhrannik}



\begin{quote}
Before the new one: say the Russian for a cashier, then the Russian for a hairdresser.
\end{quote}

\subsection*{You'll want to know: \ru{охранник}}
\textbf{\ru{охранник}} (\emph{okhránnik}) — \textquotedblleft{}a security guard\textquotedblright{}.

\begin{cousinweb}
\textbf{\ru{охранник}} is built on \emph{\ru{охрана}}, \textquotedblleft{}a guard\textquotedblright{}, on \emph{\ru{хранить}}, \textquotedblleft{}to keep\textquotedblright{}, which Vasmer calls Church Slavonic, beside native \emph{\ru{хоронить}}, now \textquotedblleft{}to bury\textquotedblright{}.
\end{cousinweb}

\begin{grammarlens}[title={What you've built}]
One more word for this chapter.
\end{grammarlens}

\subsection*{Your turn}
\begin{itemize}
\raggedright
  \item \emph{Say it:} \emph{okhránnik}
  \item \emph{Say it:} \emph{okhránnik}, once more
  \item \emph{Say it:} say \emph{okhránnik}
  \item \emph{Recall:} say the Russian for a cashier, then the Russian for a hairdresser, then say \emph{okhránnik} again
\end{itemize}

\begin{tcolorbox}[breakable,colback=teal!4,colframe=teal!35!black,title={Before you move on}]
What does \ru{охранник} mean? (\textquotedblleft{}A security guard\textquotedblright{}.) Say it once more.
\end{tcolorbox}
\section[\texorpdfstring{\ru{няня} (nyánya)}{няня (nyánya)}]{\ru{няня} (nyánya) — a nanny}

\label{lesson:RU-C230-nyanya}



\begin{quote}
Before the new one: say the Russian for a hairdresser, then the Russian for a security guard.
\end{quote}

\subsection*{You'll want to know: \ru{няня}}
\textbf{\ru{няня}} (\emph{nyánya}) — \textquotedblleft{}a nanny\textquotedblright{}.

\begin{cousinweb}
\textbf{\ru{няня}} is a nursery word, says Vasmer, of uncertain age, like Greek \emph{nánnē}, \textquotedblleft{}aunt\textquotedblright{}, Sanskrit \emph{nanā}, \textquotedblleft{}mummy\textquotedblright{}, and Latin \emph{nonna}, \textquotedblleft{}a nun\textquotedblright{}.
\end{cousinweb}

\begin{grammarlens}[title={What you've built}]
One more word for this chapter.
\end{grammarlens}

\subsection*{Your turn}
\begin{itemize}
\raggedright
  \item \emph{Say it:} \emph{nyánya}
  \item \emph{Say it:} \emph{nyánya}, once more
  \item \emph{Say it:} say \emph{nyánya}
  \item \emph{Recall:} say the Russian for a hairdresser, then the Russian for a security guard, then say \emph{nyánya} again
\end{itemize}

\begin{tcolorbox}[breakable,colback=teal!4,colframe=teal!35!black,title={Before you move on}]
What does \ru{няня} mean? (\textquotedblleft{}A nanny\textquotedblright{}.) Say it once more.
\end{tcolorbox}
